% TOP_PROBLEM: {"problemId": "problem.greek-letter-redshift-conjecture-in-algebraic-k-theory", "canonicalTitle": "Greek-Letter Redshift Conjecture in Algebraic K-Theory", "releaseRank": 453, "primaryDomain": "algebra_representation_category", "primaryDomainLabel": "Algebra, representation & category theory", "displayStatus": "open_with_solved_subcases", "releaseStatus": "open_with_solved_subcases"}
% !TeX program = lualatex
% ATTEMPT_STATUS: unresolved
% Model: GPT-6 Astra Ultra; continuation dated 27 September 2026.
\documentclass[11pt]{article}
\usepackage[margin=25mm]{geometry}
\usepackage{fontspec}
\setmainfont{Latin Modern Roman}
\usepackage{amsmath,amssymb,mathtools}
\usepackage{unicode-math}
\setmathfont{Latin Modern Math}
\usepackage{xurl,hyperref}
\hypersetup{colorlinks=true,urlcolor=blue}
\setcounter{secnumdepth}{0}
\setlength{\parindent}{0pt}
\setlength{\parskip}{5pt}
\setlength{\emergencystretch}{3em}
\begin{document}
\section*{\texorpdfstring{Greek-Letter Redshift Conjecture in Algebraic K-Theory}{Greek-Letter Redshift Conjecture in Algebraic K-Theory}}\label{453}
\subsection{Definitions and mathematical statement}
Fix a prime at which the $n$-th and $(n+1)$-st Greek-letter families in the source exist as nonzero stable homotopy classes. Let $S$ be the sphere spectrum, so $\pi_jS$ is its $j$-th stable homotopy group. Greek-letter elements are the periodic families constructed from powers of chromatic periodicity classes and connecting maps, or equivalently the specified permanent cycles in the chromatic and $BP$-Adams spectral sequences; use Section 1 of Angelini-Knoll for the indexing.

A commutative ring spectrum $R$ detects a family $\{\alpha_k^{(n)}\}$ when every member remains nonzero under the unit:
\[
 \pi_*S\longrightarrow\pi_*R.
\]
Writing $K(R)$ for the algebraic K-theory spectrum of perfect $R$-modules, Conjecture 1.1 asks that detection of the $n$-th family by $R$ imply detection of the $(n+1)$-st by $K(R)$. The source's existence qualification on the families matters, especially at small primes. Detecting a single element or merely proving a height bound is not the entire assertion.
\par\smallskip\textit{Formulation and concept references:} \href{https://arxiv.org/abs/1810.10088}{[S1]}.

\subsection{Short English statement}
Does taking algebraic K-theory raise by one the Greek-letter family detected by a commutative ring spectrum?
\subsection{Research attempt}

\paragraph{Scope and the distinction from qualitative redshift.}
For a ring spectrum $E$, write $\eta_E:S\to E$ for its unit and
$h_E:\pi_*S\to\pi_*E$ for the induced map. The target is nonvanishing
of every specified existing member of the next Greek-letter family
under $h_{K(R)}$. Existence of some nonzero homotopy class in the same
degree, periodicity in a homology theory, and chromatic-height growth
are different assertions.

The source audit on 27 September 2026 found two especially relevant
primary statements. Angelini-Knoll [S1, Theorem 1.2] proves detection
of the $\beta_k$ with $k\not\equiv0\pmod p$ in
$K(K(\mathbb F_q)_p)$, for $p\geq5$ and the specified prime-power
generator $q$ modulo $p^2$. The surrounding discussion explicitly
restricts this argument to classes surviving reduction to $V(1)$.
It does not establish the entire printed conjecture for arbitrary $R$.
Burklund--Schlank--Yuan [E1, Theorem E] prove a redshift theorem for
chromatic height of commutative ring spectra. That theorem does not
state the required elementwise unit-detection assertion. Neither
result is enlarged in scope below.

This attempt develops exact functorial reductions, proves a split-unit
subcase covering augmented sphere algebras, and audits the coefficient
reductions that occur in the known height-two example. It does not
claim to reproduce the deep spectral-sequence computation of [S1].

\paragraph{Functoriality has a useful but asymmetric direction.}
\textbf{Lemma 453.1.} Let $f:E\to F$ be a unital map of ring spectra.
For any $\alpha\in\pi_*S$, if $h_F(\alpha)\ne0$ then
$h_E(\alpha)\ne0$.

\emph{Proof.}
Unitality says $f\eta_E=\eta_F$, so
$\pi_*f(h_E(\alpha))=h_F(\alpha)$. A zero element cannot have a
nonzero image.\hfill$\square$

Thus detection in a target transfers backwards along a unital map.
It does not generally transfer forwards, since $\pi_*f$ may kill the
class. For example the unital map $S\to H\mathbb Q$ kills every
positive-degree torsion sphere class, although the sphere detects it.

For a map of commutative ring spectra $R\to R'$, extension of scalars
sends perfect $R$-modules to perfect $R'$-modules and gives
$K(R)\to K(R')$. Perfectness is preserved because finite cell modules
and their retracts remain so after extension of scalars. The map is
compatible with the units, represented by the free rank-one modules.
Consequently, proving detection in $K(R')$ suffices to prove detection
in $K(R)$. This can be used element by element: the target $R'$ may
depend on which class is being tested, provided an actual unital map
and a valid detection theorem for that target are supplied each time.

The hypothesis of the conjecture does not itself supply such a map to
a known computable detector. Reversing the implication of Lemma 453.1,
or inventing a map from a mere chromatic-height comparison, would be
an invalid reduction.

\paragraph{The trace proves preservation of already detected classes.}
We use the standard natural Dennis trace [E3] and the augmentation of
topological Hochschild homology for a commutative ring spectrum:
\[
 K(R)\stackrel{\operatorname{tr}}{\longrightarrow}\operatorname{THH}(R)
       \stackrel{\epsilon}{\longrightarrow}R.
 \tag{453.1}
\]
The augmentation is induced by collapsing the circle in
$\operatorname{THH}(R)\simeq R\otimes S^1$, where the tensor is
in commutative ring spectra. These standard trace maps are compatible
with units. Equivalently, the free rank-one module maps to the class
of $1\in\pi_0R$, so the composite of (453.1) with
$\eta_{K(R)}$ is $\eta_R$. Maps from $S$ to a spectrum are classified
by $\pi_0$, which also verifies this equality of unit maps.

\textbf{Proposition 453.2.} If $R$ detects a collection of sphere
classes, then $K(R)$ detects the same collection.

\emph{Proof.}
If $h_{K(R)}(\alpha)=0$, applying (453.1) gives
$h_R(\alpha)=0$. Taking the contrapositive proves the assertion.
\hfill$\square$

This proves preservation, not a shift from the $n$-th to the
$(n+1)$-st family. Indeed, the map from $R$ into
$\operatorname{THH}(R)$ also preserves units. Any sphere class killed
by the unit of $R$ is therefore killed by the unit of
$\operatorname{THH}(R)$. A proposed argument detecting new classes
only in the underlying $\operatorname{THH}(R)$ would miss this
obstruction. Cyclotomic or homotopy-fixed-point refinements used in
trace calculations carry additional information; they must not be
silently replaced by the underlying Hochschild-homology spectrum.

\paragraph{A complete split-unit subcase in every allowed height.}
\textbf{Proposition 453.3.} Suppose $R$ admits a unital augmentation
$a:R\to S$. Then the unit $S\to K(R)$ has a retraction as a map
of spectra. Hence $K(R)$ detects every nonzero stable sphere class,
and in particular every existing Greek-letter family in the statement.

\emph{Proof.}
Compose the map $K(a):K(R)\to K(S)$ with the Dennis trace
\[
 K(S)\longrightarrow\operatorname{THH}(S)\simeq S.
\]
Since $\operatorname{THH}(S)$ is the sphere and both maps preserve
units, the composite with $S\to K(R)$ is the identity of $S$.
Therefore the induced homomorphism $\pi_*S\to\pi_*K(R)$ is split
injective.\hfill$\square$

Here the input $R$ also detects every sphere class, since
$a\eta_R=\operatorname{id}_S$. Thus this is a genuine class of
instances satisfying the hypothesis and conclusion of the conjecture,
not a vacuous implication. It includes $R=S$ itself and commutative
spherical monoid algebras $\Sigma^\infty_+M$ with the augmentation
induced by $M\to *$, for commutative monoids $M$. All arguments can
also be made after localization at the fixed prime, replacing the
sphere consistently by $S_{(p)}$.

A unit map into an arbitrary $R$ need not split, so the proposition
does not extend to every ring detecting one selected family. Detection
of that family asserts injectivity on its individual chosen nonzero
members, not a splitting of the full unit as a spectrum map.

Finite products offer another exact propagation rule. Perfect modules
over $R_1\times R_2$ split using the two central idempotents, giving
\[
 K(R_1\times R_2)\simeq K(R_1)\times K(R_2).
\]
The sphere unit maps diagonally. A class detected in either factor is
detected in the product. If neither factor detects a class, its product
image is zero. This uses the actual component maps, not a maximum of
chromatic heights as a substitute for detecting a specified element.

\paragraph{Chromatic height cannot determine the unit image.}
The following elementary obstruction is useful when interpreting [E1].
If all homotopy groups of a ring spectrum $E$ are torsion-free, then
every positive-degree torsion sphere class maps to zero in $\pi_*E$:
its image has order dividing its finite order and must vanish.

For example, the usual Lubin--Tate spectra $E_h$ are commutative ring
spectra of chromatic height $h$. Their coefficient groups are even
periodic with coefficient ring
\[
 \pi_*E_h\cong W(k)[[u_1,\ldots,u_{h-1}]][u^{\pm1}],\qquad |u|=2
\]
up to the harmless choice of the sign of the periodicity degree.
For perfect fields $k$ of characteristic $p$, these groups are
torsion-free as abelian groups, and odd groups vanish. Thus their
units kill positive-degree torsion sphere classes, despite their
nonzero height. This uses the standard coefficient description of
Lubin--Tate theory, also part of the setting of [E1].

It follows that a numerical height statement cannot by itself imply
detection of a specified positive-degree Greek-letter class. This is
not a counterexample to the conjecture: these examples do not satisfy
the required family-detection hypothesis. It identifies the missing
information in an attempted deduction solely from qualitative redshift.
Moreover, periodicity is usually measured after smashing with a finite
complex. A nonnilpotent self-map on that complex is not the same
object as the torsion sphere class produced by connecting maps.

\paragraph{A low-height arithmetic calculation, with the unit input separated.}
Let $p$ be odd and let the prime power $q$ generate
$(\mathbb Z/p^2\mathbb Z)^\times$. Quillen's finite-field calculation,
used in [S1], gives
\[
 K_{2m-1}(\mathbb F_q)\cong\mathbb Z/(q^m-1),\qquad
 K_{2m}(\mathbb F_q)=0\quad(m\geq1).
\]
At $m=(p-1)k$, the $p$-primary cyclic group has order
$p^{1+\nu_p(k)}$. Here is a direct verification of the exponent.

Since $q$ has order $p(p-1)$ modulo $p^2$, write
$q^{p-1}=1+pc$ with $p\nmid c$. If $B=1+p^rb$ with $r\geq1$
and $p\nmid b$, binomial expansion for odd $p$ shows
\[
 \nu_p(B^p-1)=r+1.
\]
The linear term has valuation $r+1$, every intermediate term has
valuation at least $1+2r\geq r+2$, and the final term has valuation
$pr\geq r+2$. Raising to a positive exponent coprime to $p$ leaves
the valuation equal to $r$, again by the binomial expansion.
Writing $k=p^su$ with $p\nmid u$ proves
\[
 \nu_p(q^{(p-1)k}-1)=1+s=1+\nu_p(k).
 \tag{453.2}
\]

This calculation locates nonzero target groups. It does not determine
the unit homomorphism from the sphere. The separate image-of-$J$
identification discussed in [S1] is the established homotopy-theoretic
input showing that the appropriate $\alpha$-family classes have
nonzero unit images in $K(\mathbb F_q)_p$. Without that input, the
displayed group orders alone would not prove detection: a map to a
nonzero group can be the zero map.

\paragraph{The degrees in the height-two construction.}
At $p\geq5$, let $V(0)=S/p$ and choose the standard periodic map
$v_1:\Sigma^{d_1}V(0)\to V(0)$ with $d_1=2p-2$.
Its cofiber is $V(1)$. A $v_2$ self-map of $V(1)$ has degree
$d_2=2p^2-2$ at these primes. Write $i:S\to V(1)$ for the bottom
cell map and $q_\mathrm{top}:V(1)\to\Sigma^{d_1+2}S$ for the
composite of the two cofiber boundary maps to the top cell.
The basic $\beta$-family construction is represented by
\[
 \Sigma^{kd_2}S\stackrel{\Sigma^{kd_2}i}{\longrightarrow}
 \Sigma^{kd_2}V(1)\stackrel{v_2^k}{\longrightarrow}V(1)
 \stackrel{q_\mathrm{top}}{\longrightarrow}\Sigma^{d_1+2}S.
\]
After desuspension its degree is
\[
 |\beta_k|=k(2p^2-2)-2p.
 \tag{453.3}
\]
This is consistent with the indexing in [S1]. Higher families and
small primes require their own existence and indexing conventions;
we do not extrapolate (453.3) as a universal formula for them.

Smashing this diagram with $K(R)$ produces the corresponding unit
image after beginning with the unit. Nonzero periodic classes somewhere
in $\pi_*(V(1)\wedge K(R))$ do not automatically make this composite
nonzero: the bottom-cell class, the chosen periodic map and the
top-cell connecting maps all matter. The connecting map can kill a
nonzero intermediate class. No commutative ring structure on $V(1)$
is assumed or needed here, and $K(V(1)\wedge R)$ is not substituted
for the coefficient spectrum $V(1)\wedge K(R)$.

\paragraph{Reduction modulo periodicity detects only in one direction.}
For any spectrum $E$, the map $i\wedge1:E\to V(1)\wedge E$
has the elementary implication
\[
 (i\wedge1)_*(x)\ne0\quad\Longrightarrow\quad x\ne0.
 \tag{453.4}
\]
Thus a nonzero reduced image of a sphere class gives a genuine
nonzero image in $E$. The converse need not hold. Already the first
cofiber sequence
\[
 E\stackrel{p}{\longrightarrow}E\longrightarrow V(0)\wedge E
\]
shows, by exactness, that the kernel of
$\pi_dE\to\pi_d(V(0)\wedge E)$ is $p\pi_dE$.
Even a nonzero element of order $p$ can lie in this kernel: the
element $p$ in $\mathbb Z/p^2$ has precisely that property.

The second cofiber sequence gives
\[
 \ker\bigl(\pi_d(V(0)\wedge E)\to\pi_d(V(1)\wedge E)\bigr)
 =\operatorname{im}\bigl(v_1:\pi_{d-d_1}(V(0)\wedge E)
                              \to\pi_d(V(0)\wedge E)\bigr).
 \tag{453.5}
\]
Hence a class can survive the mod-$p$ step and still disappear under
the next reduction because it is $v_1$-divisible. This is why the
qualification in [S1] matters: its proof detects $\beta_k$ with
$p\nmid k$, whereas the indicated $\beta_{pk}$ are invisible to
the reduction used there. Their disappearance in that reduction
does not prove that their original unit images in iterated K-theory
are zero.

\paragraph{An exact certificate that a trace approach would need.}
Fix one existing next-family class $\beta\in\pi_dS$.
It is sufficient to construct a spectrum map $\tau:K(R)\to D$ and
prove
\[
 (\tau\eta_{K(R)})_*(\beta)\ne0.
 \tag{453.6}
\]
This immediately forces $h_{K(R)}(\beta)\ne0$, regardless of whether
$\tau$ is injective on other classes. A cyclotomic trace followed by
a suitable coefficient or fixed-point comparison is one possible
source of such a map. What must be verified is the actual composite
on the sphere, not merely the presence of a class in $\pi_dD$.

If a spectral sequence computes $\pi_*D$, finding a named class on
its first page is insufficient. One must identify it as the image of
the specified sphere class, exclude outgoing differentials and incoming
boundaries, and use a convergence statement making its nonzero
associated-graded class a nonzero element of the abutment. For a
strongly convergent separated filtration, the last implication is
formal; the identification and survival are the substantive work.
Proving (453.6) for finitely many $k$ does not establish all-family
detection without a uniform argument covering the remaining indices.

\paragraph{Diagnostics, outcome and remaining gap.}
The exact arithmetic checks accompanying this attempt verify (453.2)
for explicit prime-power generators and a range of indices, and check
the degree shifts in (453.3). They are checks of arithmetic and indexing,
not computations of stable stems or spectral-sequence differentials.
The nontrivial detection theorem for iterated finite-field K-theory is
attributed to [S1], with its restriction on indices retained.

The proved results are functorial backward propagation, preservation
of already detected classes via the trace, and full detection whenever
$R$ admits a sphere augmentation. The exact-sequence analysis explains
both the usefulness and the limitations of mod-$(p,v_1)$ detection.
The first unresolved step for arbitrary $R$ is constructing and proving
nonzero composites of the form (453.6) for every existing member of the
next family from only the given previous-family detection hypothesis.
The known height theorem and the calculations above do not provide
that missing implication. The printed conjecture remains unresolved
by this attempt.
\subsection{Sources}
\begin{itemize}
\item[S1]G. Angelini-Knoll, Detecting beta elements in iterated algebraic K-theory, Trans. AMS.
\url{https://arxiv.org/abs/1810.10088}.
\textit{Formulation source. Consulted 24 September 2026: Conjecture 1.1 and the family-detection definition in its full-text HTML version.}
\item[E1]Robert Burklund, Tomer M. Schlank and Allen Yuan,
\emph{The Chromatic Nullstellensatz}, arXiv:2207.09929.
\url{https://arxiv.org/abs/2207.09929};
\url{https://arxiv.org/html/2207.09929}.
\textit{Primary Theorem E and the chromatic-height formulation
inspected 27 September 2026. Height growth is not substituted for
unit detection of individual Greek-letter classes.}
\item[E2]Gabriel Angelini-Knoll, full-text version of [S1],
\url{https://arxiv.org/html/1810.10088}.
\textit{Conjecture 1.1, Theorem 1.2 and the preceding mod-$(p,v_1)$
qualification inspected 27 September 2026. The condition $p\geq5$,
the prime-power generator hypothesis and the restriction $p\nmid k$
on the detected basic beta elements are retained.}
\item[E3]Andrew J. Blumberg, David Gepner and Gon\c calo Tabuada,
\emph{Uniqueness of the multiplicative cyclotomic trace}, Advances
in Mathematics 260 (2014), 191--232.
\url{https://arxiv.org/abs/1103.3923};
\url{https://doi.org/10.1016/j.aim.2014.02.004}.
\textit{Primary abstract inspected 27 September 2026 for the natural
multiplicative Dennis trace and multiplicative K-theory input.}
\end{itemize}
\end{document}
